\chapter{Normal Variants and Anatomic Mimics}

\label{ch:ch31_normal_variants}

\section{Definition}
A normal variant is a benign anatomic, developmental, physiologic, or technical appearance that may mimic disease. Variants can be common, uncommon, or clinically important in a particular context.

\subsection{Why Variants Matter}
Normal variants cause two problems: they may be \textbf{mistaken for pathology} (leading to unnecessary investigation or anxiety), or they may \textbf{mask genuine disease} (because the reader assumes the finding is a known variant). Recognising variants requires knowledge of anatomy, developmental biology, and the imaging characteristics that distinguish benign from malignant.

\section{Common Normal Variants}

\begin{table}[h]
\centering
\caption{Common normal variants by body region.}
\label{tab:normal_variants}
\begin{tabular}{p{3.5cm} p{3.5cm} p{7cm}}
\toprule
\textbf{Region} & \textbf{Variant} & \textbf{Key Features} \\
\midrule
Brain & Falx cerebri calcification & Linear or nodular calcification along the interhemispheric fissure; increases with age \\
Brain & Choroid plexus calcification & Often bilateral and symmetric in the glomus of the lateral ventricles; commonly physiologic in adults and more frequent with age \\
Brain & Prominent arachnoid granulations (Pacchioni) & Smooth, well-defined, CSF-intensity lesions along the superior sagittal sinus; may erode the inner table of the skull \\
Chest & Dextrocardia & Mirror-image cardiac position; situs inversus if abdominal organs are also reversed \\
Kidney & Dromedary hump & Focal contour prominence of the left kidney due to splenic impression; shows normal enhancement \\
Liver & Prominent caudate lobe & May be a normal anatomic configuration, but caudate enlargement can also occur with chronic liver disease; assess the whole liver and clinical context \\
Musculoskeletal & Accessory ossification centers & Persistent unfused centers (e.g.\ os acromiale, bipartite patella) with smooth, well-corticated margins; bilaterality varies by variant \\
Spleen & Accessory spleen & Usually a small, well-defined nodule near the splenic hilum or pancreatic tail that typically follows splenic signal and enhancement \\
\bottomrule
\end{tabular}
\end{table}

\begin{pearlbox}[Clinical Pearl: Accessory spleen vs. recurrence]
After splenectomy for selected haematologic conditions, residual accessory splenic tissue or splenosis may contribute to persistent or recurrent disease. When clinically relevant, assess the splenic bed and consider possible ectopic sites. Such tissue often resembles the spleen in signal and enhancement, but size and appearance vary; indeterminate lesions may need confirmatory imaging.
\end{pearlbox}

\section{Safe Recognition}
Use age, symptoms, prior imaging, distribution, enhancement, and stability. Do not label an appearance as a variant solely because it is unfamiliar. If a lesion is indeterminate or has suspicious features, recommend appropriate further evaluation rather than reassurance.

\subsection{Distinguishing Variants from Pathology}
The following features favour a normal variant:
\begin{itemize}
    \item \textbf{Characteristic location} -- a dromedary hump is typically a contour bulge along the lateral midportion of the left kidney, related to splenic impression.
    \item \textbf{Characteristic morphology} -- accessory ossification centers typically have smooth, well-corticated margins at an expected ossification site; correlate with age, symptoms and trauma history.
    \item \textbf{Relationship to parent tissue} -- an accessory spleen usually has signal and enhancement similar to the spleen, although the match may vary with size, phase and technique.
    \item \textbf{Stability} -- long-term stability can support a benign variant, but stability alone does not prove that a finding is benign.
    \item \textbf{Bilaterality} -- bilateral findings can support some variants, but bilaterality alone does not establish that a finding is benign.
\end{itemize}

\subsection{When to Recommend Follow-Up}
If a finding does not meet the criteria above, do not dismiss it as a variant. Recommend follow-up imaging or further characterisation when:
\begin{itemize}
    \item The patient has symptoms or risk factors for malignancy.
    \item The morphology is atypical for a known variant.
    \item There is no prior imaging for comparison.
    \item The finding has changed in size or appearance.
\end{itemize}

\begin{pearlbox}[Clinical Pearl: Os acromiale]
Os acromiale is an unfused ossification center of the acromion that can be confused with an acromial fracture. Smooth, well-corticated margins at a typical location support a variant; assess for acute injury and correlate with symptoms. Bilaterality may support the diagnosis but is not required.
\end{pearlbox}

\section{Communication}
Explain the distinguishing features and any relevant caveat. Avoid implying that every variant is harmless in every patient; some variants alter procedural risk or are associated with disease in selected settings.

\subsection{Structured Reporting of a Suspected Variant}
When reporting a finding that you believe to be a normal variant:
\begin{enumerate}
    \item State the finding clearly and provide measurements.
    \item Describe the features that support it being a variant (location, morphology, signal characteristics).
    \item State whether prior imaging is available for comparison.
    \item If confident, state ``consistent with normal variant'' and recommend no further follow-up.
    \item If uncertain, recommend targeted follow-up and state the specific concern.
\end{enumerate}

\begin{pearlbox}[Impression:]
Many normal variants have no clinical consequence, but some affect procedural planning or are relevant in particular diseases. Recognition should use the full pattern of location, morphology, signal or attenuation, enhancement, prior imaging and clinical context; no single feature proves a finding is benign. When uncertain, recommend appropriate further evaluation rather than reassurance.
\end{pearlbox}

\clearpage
